% !TeX program = pdflatex
\documentclass[10pt,a4paper]{article}

\usepackage[utf8]{inputenc}
\usepackage[T1]{fontenc}
\usepackage{geometry}
\usepackage[hidelinks]{hyperref}
\usepackage{xcolor}

\geometry{
  top=1.7cm,
  bottom=1.8cm,
  left=1.8cm,
  right=1.8cm
}

\pagestyle{empty}
\frenchspacing
\setlength{\parindent}{0pt}
\setlength{\parskip}{0.45em}

\definecolor{mydarkgreen}{RGB}{0,0,0}

\begin{document}

\begin{minipage}[t]{0.48\textwidth}
\textbf{\large Javier Andres TARAZONA JIMENEZ}\\
Étudiant ingénieur IA \\
Machine Learning, Vision 3D, LLMs et prototypage \\
Massy, Île-de-France, France\\
\href{mailto:javitar06@gmail.com}{javitar06@gmail.com}\\
+33 07 46 36 97 75\\
\href{https://www.linkedin.com/in/javier-andres-tarazona-jimenez-84b489222/}{LinkedIn}
\end{minipage}
\hfill
\begin{minipage}[t]{0.42\textwidth}
\raggedleft
3DS Talent Acquisition\\
Dassault Systèmes\\
Vélizy-Villacoublay, Île-de-France, France\\
\end{minipage}

\vspace{1.1cm}

\raggedleft
Massy, le 19 juin 2026

\vspace{0.7cm}

\raggedright
\textbf{Objet : Candidature à l'alternance}

\vspace{0.5cm}

{\small

Madame, Monsieur,

Actuellement étudiant colombien en double diplôme entre l’Universidad Nacional de Colombia et ENSTA Paris, au sein de l’Institut Polytechnique de Paris, je souhaite vous présenter ma candidature pour une alternance en intelligence artificielle chez Dassault Systèmes à partir de septembre 2026.

J’ai choisi l’informatique par intérêt pour la compréhension des systèmes complexes, par goût pour la construction de solutions concrètes à partir d’idées abstraites, et par volonté d’acquérir une autonomie technique solide. L’intelligence artificielle s’est ensuite imposée comme une évolution naturelle de ce parcours : passer de la programmation de systèmes à la conception de systèmes capables de percevoir, d’apprendre et d’aider à la décision. Ce domaine m’attire parce qu’il combine rigueur mathématique, incertitude, expérimentation et problèmes ouverts. La vision par ordinateur, en particulier, relie fortement mon côté appliqué à ma curiosité scientifique : elle transforme les mathématiques et le code en perception exploitable du monde réel.

L’intérêt que je porte à Dassault Systèmes dépasse largement l’image d’une entreprise de logiciels 3D. Je vois surtout en vous un acteur majeur de la simulation, des jumeaux virtuels, de la géométrie numérique et des plateformes industrielles. Votre valeur réside dans la capacité à relier l’IA, les données, la géométrie, la physique, l’optimisation et l’ingénierie pour améliorer la conception et la prise de décision dans des secteurs industriels variés. Pour un profil comme le mien, intéressé par une IA ayant un impact économique et technologique réel, cette approche est plus stimulante qu’une application isolée. Elle correspond aussi à ma manière de réfléchir : comprendre les interactions entre les composants d’un système, modéliser leurs contraintes et construire des solutions utiles dans un environnement complexe.

Vos offres en apprentissage présentent un noyau commun qui correspond précisément à mon parcours : IA appliquée à la 3D, segmentation ou reconnaissance de formes, maillages, données géométriques, prototypage, expérimentation mesurée et intégration dans des outils industriels. Mon stage actuel à ENSTA Paris – U2IS me permet de travailler sur la vision 3D auto-supervisée, la synthèse de nouvelles vues, la génération de données 3D synthétiques avec Blender, les rendus multi-vues, les métadonnées de caméras et le pré-entraînement avec PyTorch et Vision Transformers. Cette expérience m’a appris à construire des pipelines expérimentaux, préparer des données, comparer des configurations et documenter des résultats de manière rigoureuse.

Je peux également apporter une base solide en développement logiciel et en vision par ordinateur. Chez Engin A.I, j’ai travaillé sur l’analyse d’images, l’optimisation de pipelines OpenCV/NumPy et la structuration de bibliothèques Python en architecture modulaire. Dans mes projets académiques, j’ai développé des systèmes de segmentation, de détection et de suivi visuel, ainsi que des applications full-stack. Cette double compétence <<IA appliquée et ingénierie logicielle>> me semble essentielle pour transformer une idée de recherche en prototype fonctionnel, testable et progressivement industrialisable.

Rejoindre Dassault Systèmes représenterait pour moi une transition structurante entre l’IA académique et l’IA industrielle. J’y vois l’opportunité d’apprendre comment des modèles, des algorithmes et des prototypes deviennent des composants intelligents intégrés à des plateformes utilisées par de grandes industries. Ce cadre me permettrait de faire évoluer mon profil de spécialiste en vision 3D vers un profil plus complet : IA, 3D, simulation, géométrie, ingénierie système et produit réel.

Je serais heureux de pouvoir échanger avec vous afin de vous présenter plus en détail ma motivation et l’adéquation de mon parcours avec vos besoins.

\textbf{Cordialement},

Javier Andres TARAZONA JIMENEZ


}

\end{document}